\documentclass[11pt]{article}
\usepackage[a4paper,margin=28mm]{geometry}
\usepackage{microtype}
\usepackage[hidelinks]{hyperref}
\title{A Sample Review}
\author{Example Project}
\date{}

\begin{document}
\maketitle

\begin{abstract}
This document demonstrates a local review workflow for a LaTeX draft.
Comments remain beside the source while the compiled PDF provides context.
All text and numbers in this example are fictional.
\end{abstract}

\section{A clear starting point}
A useful introduction tells the reader what question the document will answer.
It gives enough context to understand the question before introducing details.
The opening paragraph should make the purpose of the comparison explicit.

\section{Comparing two versions}
We compare an initial draft with a revised draft of the same short document.
The revision changes the explanation while keeping the example conditions fixed.
This makes it easier to discuss whether the new wording answers the question.

\subsection{Review procedure}
First, read the passage and select the text you want to discuss.
Then add a comment, check the PDF, and reply when the wording is clearer.
Use a saved checkpoint to inspect what changed between review rounds.

\begin{table}[ht]
\centering
\caption{Fictional values used only to demonstrate a review.}
\label{tab:example}
\begin{tabular}{lrr}
\hline
Draft & Reading time (s) & Questions raised \\
\hline
Initial & 45 & 4 \\
Revised & 38 & 2 \\
\hline
\end{tabular}
\end{table}

\section{Interpreting the example}
Table~\ref{tab:example} supplies a small amount of numerical context.
These fictional values are not evidence that one writing strategy is better.
A real study would need a defined protocol and observations from readers.

\newpage
\section{Finishing the review}
Check that each reply explains the change made to the corresponding passage.
Resolve a thread after confirming the revision, or reopen it for more discussion.
Export a PDF with comments when another reader needs to review the conversation.

\end{document}
